\begin{table}[htbp]
   
   \caption{\label{tab:6yr_huc02fe_inorg_ravalli_2005_k2_elas_prod} Effect of upstream coal production on inorganic chemicals by upstream distance, elasticities at the sample means, SYR2 1998--2005}
   \bigskip
   
   \centering
   
\begin{adjustbox}{width = \textwidth, center}
   \begin{tabular}{lcccc}
      \toprule
       & Arsenic & Nitrate & Barium & Selenium \\
                                         & (1)           & (2)         & (3)     & (4)\\  
      \midrule 
      \multicolumn{5}{l}{\textbf{Panel A: Within one HUC12 upstream}} \\
      Upstream coal prod. (rel. to mean) & 0.096$^{***}$ & 0.022$^{*}$ & 0.010   & 0.038\\   
                                         & (0.029)       & (0.013)     & (0.011) & (0.025)\\   
       \\
      \midrule 
      \multicolumn{5}{l}{\textbf{Panel B: Within two HUC12 upstream}} \\
      Upstream coal prod. (rel. to mean) & 0.064   & 0.045$^{***}$ & 0.089   & 0.046\\   
                                         & (0.047) & (0.011)       & (0.080) & (0.048)\\   
       \\
      Observations                       & 366     & 472           & 353     & 350\\  
      \bottomrule
   \end{tabular}
\end{adjustbox}
   
   {\tiny\linespread{1}\selectfont \par \raggedright 
   \textit{Notes:} Within each chemical, the column shows mean measured concentration from the EPA 6-Year Review. Mean concentration and the explanatory variable (coal production during the year (in millions of short tons) in watersheds within one flow step upstream of the utility's intake (Panel A) or within two flow steps upstream (Panel B)) are each divided by their sample mean within each chemical and panel, so coefficients are elasticities evaluated at the sample means. Sample: utilities with a coal mine within two flow steps upstream of their intake and no coal mine colocated with their intake. Standard errors clustered at the utility level. All specifications include utility and HUC02 $\times$ year fixed effects. *** p$<$0.01, ** p$<$0.05, * p$<$0.1.}
   
\end{table}


